\section{Projektgegenstand und Kurzfassung}

Das Projekt entwickelt ein lokal betreibbares System zur
\emph{Personen-Re-Identifikation} (Person Re-Identification, ReID). Es erkennt
Personen in Videos oder einem lokalen Kamerabild, verfolgt sie innerhalb der
Aufnahme und ordnet erneute Sichtungen einer synthetischen Personen-ID zu. Das
System beantwortet damit nicht die Frage nach einem Klarnamen, sondern die
technische Frage: \enquote{Handelt es sich mit hinreichender Evidenz um dieselbe
bereits beobachtete Person?}

Der aktuelle Prototyp verbindet YOLO-basierte Personendetektion, ByteTrack oder
BoT-SORT für Multi-Object-Tracking, Torchreid/OSNet für visuelle Embeddings,
eine mehrstufige Matching-Logik sowie SQLite oder Qdrant zur lokalen
Vektorsuche. Streamlit dient als Bedien- und Diagnoseoberfläche. Die Pipeline
speichert neben einem gemittelten Personenprofil auch einzelne hochwertige
Referenz-Embeddings und nutzt diese später erneut. Unsichere Beobachtungen
führen nicht unmittelbar zu einer endgültigen Identitätsentscheidung.

Langfristig soll die ReID-Pipeline eine lokale Kontextschicht für komplexere
Anwendungen bilden: etwa personenabhängige, aber datensparsame
Gebäudeautomation, kamerübergreifende Anwesenheitsmodelle, unterstützende
Sicherheits- und Assistenzfunktionen oder Sportanalytik. Eine direkte
IoT-Steuerung ist im heutigen Projektstand noch nicht implementiert. Das Exposé
trennt deshalb konsequent zwischen vorhandenem System, Forschungsarbeit und
Anwendungsperspektive.

\subsection{Leitfragen}

Die weitere Arbeit orientiert sich an vier Leitfragen:

\begin{enumerate}[label=\textbf{F\arabic*:},leftmargin=18mm]
  \item Wie lässt sich eine Person in Innenraumvideos lokal und ohne Klarnamen
    über zeitlich getrennte Sichtungen möglichst stabil wiedererkennen?
  \item Wie stark reduzieren Qualitätsfilter, Mehrbild-Puffer und eine
    dreistufige Evidenzentscheidung falsche Zusammenführungen und
    ID-Fragmentierung gegenüber einer einzelnen Schwellenwertentscheidung?
  \item Welche Kombination aus Detektor, Tracker, ReID-Modell und Vektorspeicher
    bietet für die Zielhardware den besten Kompromiss aus Genauigkeit,
    Nachvollziehbarkeit und Laufzeit?
  \item Wie muss eine Ereignisschnittstelle gestaltet sein, damit IoT-Systeme
    nur zweckgebundene, pseudonyme und mit Unsicherheit versehene Signale
    erhalten?
\end{enumerate}


\section{Ausgangslage, Motivation und Relevanz}

\subsection{Vom Erkennen einer Person zum Verstehen einer Situation}

Eine reine Objekterkennung kann feststellen, dass sich eine Person im Bild
befindet. Für komplexere Anwendungen reicht dies nicht aus. Soll ein System
zwischen wiederkehrenden Bewohnern, Gästen und einer unbekannten Person
unterscheiden, muss es Beobachtungen über Zeit zusammenführen. Tracking hält
eine temporäre ID innerhalb einer zusammenhängenden Sequenz; ReID soll dieselbe
Person nach Verdeckung, Wiedereintritt, einem Neustart oder später auch an einer
anderen Kamera wiederfinden. Diese Trennung ist ein zentrales
Architekturprinzip.

Im Wohn- und Gebäudeumfeld entsteht daraus ein möglicher Mehrwert: Licht,
Medien, Raumklima oder Assistenzfunktionen könnten auf wiederkehrende,
pseudonyme Kontexte reagieren. Ein IoT-System müsste dann nicht das Videobild
erhalten, sondern nur ein minimiertes Ereignis wie
\code{person\_000004 entered living\_room} einschließlich Zeitstempel,
Konfidenz und Quelle. MQTT bietet hierfür perspektivisch ein standardisiertes
Publish/Subscribe-Protokoll \citep{oasis2019mqtt}. Die Freigabe einer Aktion
darf jedoch nicht allein aus einem unsicheren KI-Match folgen; kritische
Aktionen benötigen zusätzliche Regeln oder eine zweite Bestätigung.

\subsection{Warum das Problem anspruchsvoll ist}

ReID muss mit ähnlicher Kleidung, wechselnden Outfits, Lichtänderungen,
Bewegungsunschärfe, Verdeckung, kleinen Bildausschnitten und veränderten
Perspektiven umgehen. Im Haus kommen enge Räume, Türdurchgänge, wechselnde
Kameraabstände und gleichzeitige Personen hinzu. OSNet adressiert die visuelle
Wiedererkennung durch eine leichtgewichtige, mehrskalige Merkmalsrepräsentation
\citep{zhou2019osnet}; dennoch liefert auch ein leistungsfähiger Encoder nur
Ähnlichkeiten und keine sichere Identität.

Die gesellschaftliche Relevanz ist ambivalent: Dieselbe Technik kann nützliche
Assistenz und lokale Automatisierung ermöglichen, zugleich aber Bewegungs- und
Verhaltensprofile erzeugen. Nach der Datenschutz-Grundverordnung können
biometrische Daten besondere Kategorien personenbezogener Daten darstellen,
wenn sie mittels spezieller technischer Verarbeitung zur eindeutigen
Identifizierung genutzt werden \citep{eu2016gdpr}. Auch der AI Act fasst
biometrische Identifikation funktional als Vergleich mit Referenzdaten
\citep{eu2024aiact}. Das Projekt benötigt deshalb eine dokumentierte
Rechtsgrundlage, Zweckbindung, Datenminimierung, Löschfristen und
Zugriffskontrolle, bevor es außerhalb kontrollierter Tests eingesetzt wird.


\section{Zielsetzung und Abgrenzung}

\subsection{Projektziele}

Das Vorhaben soll eine reproduzierbare und evaluierbare Baseline schaffen, die:

\begin{itemize}[leftmargin=7mm]
  \item Video- und lokale Kameraquellen verarbeitet,
  \item Personen detektiert und mit temporären Track-IDs verfolgt,
  \item aus ausreichend guten Personenausschnitten normalisierte
    ReID-Embeddings erzeugt,
  \item bekannte und potenziell neue Personen über mehrere Evidenzstufen
    unterscheidet,
  \item Profile kontrolliert aus hochwertigen Beobachtungen erweitert,
  \item Entscheidungen, Scores, Qualitätswerte und Alternativkandidaten
    nachvollziehbar protokolliert,
  \item SQLite und Qdrant hinter einer austauschbaren Store-Schnittstelle
    unterstützt und
  \item reproduzierbare Evaluationen für Tracking und ReID ermöglicht.
\end{itemize}

\subsection{Nicht-Ziele des aktuellen MVPs}

Das System stellt keine behördliche Identifizierung, Zugangskontrolle oder
Garantie einer Identität bereit. Es speichert standardmäßig keine Klarnamen und
verwendet weder Gesichts-, Stimm-, Emotions- noch Gangbilderkennung. Eine
produktive Multi-Kamera-Architektur, Browser-Remote-Kameras und die Steuerung
realer IoT-Geräte sind noch nicht umgesetzt. Ebenso sind Balltracking,
Spielfeldabbildung und vollständige Fußballstatistiken vorbereitete
Erweiterungspunkte, nicht Teil der stabilen ReID-Baseline.


\section{Systemaufbau}

\subsection{Datenfluss}

\begin{figure}[H]
  \centering
  \begin{tikzpicture}
    \node[stage] (input) at (0,0) {Video /\\Kamera};
    \node[stage] (detect) at (3.8,0) {YOLO\\Personendetektion};
    \node[stage] (track) at (7.6,0) {ByteTrack /\\BoT-SORT};
    \node[stage] (crop) at (11.4,0) {Person-Crop\\und Qualit\"at};
    \node[stage] (embed) at (11.4,-1.8) {OSNet-\\Embedding};
    \node[stage] (match) at (7.6,-1.8) {Top-K-Matching\\und Evidenz};
    \node[stage] (store) at (3.8,-1.8) {SQLite / Qdrant\\Profile + Events};
    \node[stage] (output) at (0,-1.8) {Live-Preview /\\Ergebnisvideo};

    \draw[flow] (input) -- (detect);
    \draw[flow] (detect) -- (track);
    \draw[flow] (track) -- (crop);
    \draw[flow] (crop) -- (embed);
    \draw[flow] (embed) -- (match);
    \draw[flow] (match) -- (store);
    \draw[flow] (store) -- (output);
  \end{tikzpicture}
  \caption{Datenfluss der implementierten ReID-Pipeline}
  \label{fig:architecture}
\end{figure}

Die Orchestrierung liegt in einer zentralen Pipeline-Klasse. Der Detektor und
Tracker liefern pro Frame Bounding Box, Confidence und Track-ID. Nach dem Crop
berechnet OSNet einen normalisierten Vektor. Die Suche vergleicht diesen mit
gespeicherten Mittelwerten und einzelnen Referenz-Samples. Der Store führt
Profile und protokolliert Ereignisse; Streamlit zeigt Parameter, Live-Bild,
Ergebnisse und Diagnosedaten.

\subsection{Modulare Projektstruktur}

\begin{table}[H]
  \centering
  \caption{Zentrale Module des Repositorys}
  \label{tab:modules}
  \begin{tabularx}{\textwidth}{L{36mm}Y}
    \toprule
    \textbf{Bereich} & \textbf{Verantwortung} \\
    \midrule
    \file{app/ui} & Streamlit-Oberfläche, Modus- und Parameterwahl, Kamera- und Ergebnisansicht \\
    \file{app/pipeline} & Detection, Tracking, ReID-Encoding, Matching und Orchestrierung \\
    \file{app/storage} & Gemeinsames Datenmodell sowie SQLite- und Qdrant-Implementierung \\
    \file{app/modes} & Default-, Max-Accuracy-, Football- und benutzerdefinierte Konfigurationen \\
    \file{app/evaluation} & Manifest, Event-Logging, Datenbank-Snapshots und Metriken \\
    \file{app/utils} & Crop-Qualität, Detailmerkmale, Bewegung, Kamera- und Modellverwaltung \\
    \file{scripts} & Installation, Start, Smoke-Test, Modellverwaltung und Evaluation \\
    \bottomrule
  \end{tabularx}
\end{table}

Die Modusabstraktion erlaubt unterschiedliche Genauigkeits- und
Anwendungsprofile, ohne die Kernpipeline zu duplizieren. Der
\code{default}-Modus verwendet derzeit YOLOv8n, ByteTrack, OSNet und SQLite.
Der ressourcenintensivere Modus kombiniert YOLOv8x, BoT-SORT und Qdrant. Der
Fußballmodus erbt bewusst die stabile ReID-Pipeline und stellt Erweiterungspunkte
für Teamfarbe, Ball, Homographie und Statistik bereit.


\section{Implementierte Technologien}

\begin{table}[H]
  \centering
  \small
  \caption{Technologiestack und Aufgabe im System}
  \label{tab:tech}
  \begin{tabularx}{\textwidth}{L{31mm}L{39mm}Y}
    \toprule
    \textbf{Ebene} & \textbf{Technologie} & \textbf{Aufgabe und Status} \\
    \midrule
    Laufzeit & Python 3.10.8, uv, uv.lock & Projektgebundene, reproduzierbar gesperrte Windows-Umgebung \\
    Video & OpenCV & Kamera-Scan, VideoCapture, Cropping, Qualitätsmaße und Ergebnisvideo \\
    Detection & Ultralytics YOLOv8 & Personklasse aus COCO-Modellen, lokale Modellgewichte, GPU/CPU \\
    Tracking & ByteTrack / BoT-SORT & Temporäre Objekt-IDs innerhalb einer Sequenz; beide Tracker sind konfigurierbar \citep{zhang2022bytetrack,aharon2022botsort} \\
    ReID & Torchreid / OSNet & Normalisierte visuelle Embeddings; OSNet-Varianten mit modellabhängiger Dimension \citep{zhou2019osnet} \\
    Kontext & OpenCV-Heuristiken, Motion & Weiche Attribute, Crop-Qualität, Bewegung und räumliche Plausibilität \\
    Speicherung & SQLite / Qdrant & Lokale Profile, Einzel-Samples, Ereignisse und Top-K-Vektorsuche; Qdrant ist optional für skalierbare Ähnlichkeitssuche \citep{qdrant2026vectors} \\
    Oberfläche & Streamlit, pandas & Lokale Parametrierung, Live-Preview und tabellarische Diagnose \\
    Hardware & PyTorch, CUDA 12.8 & Getesteter GPU-Pfad auf NVIDIA RTX 4070 sowie CPU-Fallback \\
    \bottomrule
  \end{tabularx}
\end{table}

Ultralytics kapselt Detektion und Tracking über eine gemeinsame API. Die
Projektimplementierung ruft den Tracker frameweise mit persistentem Zustand
auf und filtert auf die Personenklasse. ByteTrack assoziiert auch schwächere
Detektionen in einer zweiten Stufe und kann dadurch kurz verdeckte Tracks
erhalten \citep{zhang2022bytetrack}. BoT-SORT ergänzt unter anderem
Kamerabewegungskompensation und optionale Appearance-Information
\citep{aharon2022botsort}. Die Tracker-ID bleibt dennoch eine lokale
Sequenzkennung; die globale synthetische Personen-ID entsteht erst in der
eigenen ReID-Schicht.


\section{Methodische Prinzipien}

\subsection{Das Siebprinzip der Personenwiedererkennung}

Die zentrale Projektidee ist ein kaskadiertes Sieb: Jede Stufe verwirft
ungeeignete Daten oder reduziert die Sicherheit einer Entscheidung. Erst wenn
mehrere Prüfungen bestanden sind, darf eine Beobachtung ein Profil anlegen oder
verändern.

\begin{enumerate}[label=\textbf{S\arabic*:},leftmargin=18mm,itemsep=4pt]
  \item \textbf{Detektionssieb:} Nur YOLO-Detektionen der Klasse Person oberhalb
    der eingestellten Confidence werden übernommen.
  \item \textbf{Trackingsieb:} ByteTrack oder BoT-SORT verbindet passende Boxen
    zeitlich und erzeugt eine temporäre Track-ID.
  \item \textbf{Crop- und Qualitätssieb:} Zu kleine, unscharfe, schlecht
    belichtete, ungünstig proportionierte oder am Rand abgeschnittene Crops
    werden nicht für ReID verwendet.
  \item \textbf{Mehrbildsieb:} Ein neuer Track benötigt mehrere gute Frames.
    Aus ihnen wird ein kombiniertes Embedding erzeugt; der beste Crop dient als
    Snapshot.
  \item \textbf{Ähnlichkeitssieb:} Die normalisierten OSNet-Vektoren werden per
    Kosinusähnlichkeit gegen Mittelwert und hochwertige Einzel-Samples gesucht.
    Bereits im selben Frame verwendete Personen-IDs sind ausgeschlossen.
  \item \textbf{Kontextsieb:} Detailmerkmale und räumliche Bewegung wirken nur
    als kleine Re-Ranking- oder Plausibilitätssignale. Sie ersetzen OSNet nicht.
  \item \textbf{Evidenzsieb:} Strong, Weak/Pending und Low unterscheiden sichere
    Zuordnung, vorläufige Zuordnung und fehlende Evidenz. Eine neue Person wird
    erst nach mehreren niedrigen Treffern in einem Zeitfenster angelegt.
\end{enumerate}

Dieses Vorgehen reduziert zwei typische Fehler: Ein schlechter Einzel-Crop soll
weder ein bestehendes Profil verunreinigen noch vorschnell eine neue Person-ID
erzeugen. Gleichzeitig verhindert der Same-Frame-Ausschluss, dass zwei
gleichzeitig sichtbare Personen derselben synthetischen ID zugeordnet werden.

\subsection{Dreistufige Entscheidung statt Ja/Nein}

Für einen besten Match-Score $s$ gelten im Default-Modus beispielhaft drei
Bereiche:

\begin{equation}
  \operatorname{Entscheidung}(s)=
  \begin{cases}
    \text{Strong Match und Profilwachstum}, & s \geq 0{,}82 \\
    \text{Weak/Pending ohne Profilupdate}, & 0{,}68 \leq s < 0{,}82 \\
    \text{Low Match; weitere Evidenz sammeln}, & s < 0{,}68.
  \end{cases}
\end{equation}

Eine neue Person entsteht im Default-Modus erst, wenn mindestens sechs
Evidenzereignisse in einem Fenster von 30 Frames vorliegen und mindestens
80 Prozent davon höchstens den Score 0,58 erreichen. Die Werte sind
Konfigurationsparameter und keine allgemein gültigen Grenzwerte; sie müssen auf
einem festen Benchmark kalibriert werden.

\subsection{Wiederverwendung gesammelter Daten}

Das System lernt nicht durch unkontrolliertes Nachtrainieren des neuronalen
Netzes. Stattdessen wächst die Referenzdatenbasis kontrolliert:

\begin{itemize}[leftmargin=7mm]
  \item Ein Profil besitzt ein qualitätsgewichtetes mittleres Embedding.
  \item Zusätzlich werden einzelne hochwertige Embedding-Samples mit Quelle,
    Frame, Track, Qualität und Zeitstempel gespeichert.
  \item Neue Suchanfragen werden gegen beide Repräsentationen verglichen; pro
    Person zählt der beste Referenztreffer.
  \item Nur sichere Matches und ausreichend hochwertige Updates dürfen das
    Profil erweitern. Pending-Matches bleiben zunächst schreibgeschützt.
  \item Ein Kalibrierungsmodus kann ein neues Profil aus besonders guten
    Beobachtungen aufbauen oder ein vorhandenes Profil gezielt erweitern.
  \item Fixe Datenbank-Snapshots ermöglichen faire Vergleichstests; ein
    wachsender Datenbankmodus bildet dagegen den praktischen Lernverlauf ab.
\end{itemize}

Diese Trennung bewahrt alte, möglicherweise später hilfreiche Perspektiven,
macht Profilwachstum nachvollziehbar und vermeidet, dass ausschließlich ein
Mittelwert die Vielfalt einer Person verdeckt. Künftig müssen redundante oder
schlechte Samples begrenzt, versioniert und gezielt löschbar werden.

\subsection{Erklärbarkeit durch Ereignisprotokolle}

Pro ReID-Ereignis werden unter anderem Track-ID, Person-ID, Bounding Box,
Qualität, finaler Score, visueller Score, Detailanteil, Entscheidungszone und
Top-K-Kandidaten protokolliert. Dadurch lässt sich unterscheiden, ob ein Fehler
aus instabilem Tracking, schlechter Bildqualität, einer knappen
Ähnlichkeitsentscheidung oder einem problematischen Profilupdate entstand.


\section{Verworfene, ersetzte und zurückgestellte Ansätze}

{\small
\begin{longtable}{L{31mm}L{27mm}L{76mm}}
  \caption{Technologieentscheidungen mit Begründung}
  \label{tab:decisions} \\
  \toprule
  \textbf{Ansatz} & \textbf{Entscheidung} & \textbf{Begründung} \\
  \midrule
  \endfirsthead
  \toprule
  \textbf{Ansatz} & \textbf{Entscheidung} & \textbf{Begründung} \\
  \midrule
  \endhead
  LLM für visuelle Erkennung & verworfen & Die Aufgabe benötigt Detection, zeitliche Datenassoziation und ReID-Embeddings, keine Sprachmodell-Inferenz. \\
  32D-Farbhistogramm & ersetzt & Als Demo schnell, aber stark von Kleidung und Beleuchtung abhängig; durch OSNet ersetzt. Alte 32D- und neue 512D-Profile dürfen nicht gemischt werden. \\
  Einzelner harter Threshold & ersetzt & Ein binäres Ja/Nein erzeugte vorschnell neue IDs und unkontrollierte Updates; ersetzt durch Strong/Weak/Low und zeitliche Evidenz. \\
  Erster Crop eines Tracks & ersetzt & Häufig unscharf, klein oder abgeschnitten; ersetzt durch Quality Gate und Mehrbild-Puffer. \\
  SAM/SAM3 als ReID-Modell & verworfen & Segmentierung erzeugt Masken, aber nicht die benötigte Personenrepräsentation. Segmentierung bleibt als möglicher Vorverarbeitungstest offen. \\
  Qdrant als Pflichtdienst & nicht verworfen, optional & SQLite minimiert den MVP-Aufwand; Qdrant ist bereits als Local-, Memory- oder Server-Backend angebunden und für größere Datenmengen vorgesehen. \\
  Gesichts-, Stimm-, Emotions- und Gangbilderkennung & zurückgestellt & Höhere Sensitivität, zusätzliche Modelle und eigene rechtliche sowie ethische Bewertung erforderlich. \\
  WebRTC-Browserkamera & zurückgestellt & Aktuell greift OpenCV auf lokale Windows-Kameras zu; eine echte Remote-Browserkamera benötigt eine zusätzliche Architektur. \\
  Linux und AMD/ROCm & derzeit nicht unterstützt & Die gepflegte Baseline zielt auf Windows, NVIDIA/CUDA und CPU-Fallback; Portierung erst nach stabiler Baseline. \\
  \bottomrule
\end{longtable}
}

Die Tabelle unterscheidet bewusst zwischen \emph{verworfen}, \emph{ersetzt}
und \emph{zurückgestellt}. Beispielsweise ist Qdrant keine verworfene
Technologie, sondern eine implementierte Alternative; SAM ist nicht als
ReID-Kern geeignet, kann aber später bessere Vordergrund-Crops liefern.


\section{Anwendungsszenario: Haus und IoT}

\subsection{Zielbild}

Im Zielbild betreibt jede Kamera einen klar abgegrenzten Erfassungsprozess. Die
ReID-Schicht erzeugt aus Videodaten ein minimiertes Kontext-Ereignis. Eine
separate Regel- und IoT-Schicht entscheidet anschließend, ob und welche Aktion
zulässig ist.

\begin{figure}[H]
  \centering
  \begin{tikzpicture}[node distance=7mm]
    \node[stage] (cam) {Kamera};
    \node[stage,right=of cam] (reid) {lokale ReID-\\Pipeline};
    \node[stage,right=of reid] (event) {pseudonymes\\Kontext-Ereignis};
    \node[stage,below=of event] (policy) {Policy Engine\\und Freigabe};
    \node[stage,left=of policy] (mqtt) {MQTT / lokaler\\Ereignisbus};
    \node[stage,left=of mqtt] (iot) {IoT-Aktor\\oder Assistenz};
    \draw[flow] (cam) -- (reid);
    \draw[flow] (reid) -- (event);
    \draw[flow] (event) -- (policy);
    \draw[flow] (policy) -- (mqtt);
    \draw[flow] (mqtt) -- (iot);
  \end{tikzpicture}
  \caption{Vorgeschlagene Trennung zwischen Wahrnehmung, Entscheidung und Aktion}
  \label{fig:iot}
\end{figure}

Mögliche, zunächst nicht sicherheitskritische Demonstratoren sind
personenabhängige Lichtszenen, lokale Komfortprofile oder eine anonyme
Belegungsstatistik. Ein Ereignis sollte mindestens
\code{pseudonymous\_person\_id}, \code{zone\_id}, Zeitstempel, Konfidenz,
Entscheidungszone, Kameraquelle und eine kurze Gültigkeitsdauer enthalten.
Videoframes und Embeddings werden nicht an den IoT-Aktor weitergereicht.

\subsection{Regeln für sichere Automatisierung}

\begin{itemize}[leftmargin=7mm]
  \item Wahrnehmung, Policy und Geräteaktion bleiben getrennte Komponenten.
  \item Weak- oder Pending-Matches dürfen keine irreversible oder
    sicherheitskritische Aktion auslösen.
  \item Gleichzeitige Sichtbarkeit, Raumübergänge und zeitliche Plausibilität
    werden vor einer Aktion geprüft.
  \item Bewohner müssen Einblick, Korrektur- und Löschmöglichkeiten erhalten.
  \item Fällt ReID aus, bleibt ein sicherer manueller oder nicht-personalisierter
    Betriebsmodus verfügbar.
\end{itemize}


\section{Evaluation und erwartete Ergebnisse}

\subsection{Versuchsdesign}

Die vorhandene Evaluationsschicht verarbeitet manifestbasierte Videos. Eine
Kalibrierungsaufnahme erzeugt einen reproduzierbaren Datenbankstand. Für
\code{fixed\_db} startet jedes Testvideo mit derselben Kopie; für
\code{learn\_through} wächst das Profil über die Reihenfolge der Videos. Ein
adaptiver Kalibrierungsmodus erlaubt kontrolliertes Profilwachstum.

Die Datensätze sollen systematisch variieren:

\begin{itemize}[leftmargin=7mm]
  \item Ansicht und Pose: vorne, seitlich, hinten, sitzend, gehend;
  \item Erscheinung: Basis, Brille, anderes Oberteil und Kombinationen;
  \item Aufnahme: Licht, Entfernung, Auflösung, Bewegung und Verdeckung;
  \item Szene: Einzelperson, mehrere Personen, ähnliche Kleidung,
    Wiedereintritt sowie später unterschiedliche Kameras.
\end{itemize}

\subsection{Metriken und Vergleiche}

Bereits vorhanden sind unter anderem Anzahl und Dominanz von Track- und
Personen-IDs, Track- und Person-Switches, Profilfragmentierung sowie Strong-,
Weak- und Low-Ereignisse. Ergänzt werden sollen Precision, Recall, False Match
Rate, False Non-Match Rate sowie etablierte Trackingmetriken wie IDF1 und HOTA.

Verglichen werden mindestens:

\begin{enumerate}[leftmargin=8mm]
  \item ByteTrack gegen BoT-SORT unter identischen Bedingungen,
  \item verschiedene YOLO- und OSNet-Größen,
  \item feste Schwelle gegen das vollständige Siebprinzip,
  \item SQLite gegen Qdrant hinsichtlich Ergebnisgleichheit und Laufzeit sowie
  \item fixe gegen kontrolliert wachsende Referenzdaten.
\end{enumerate}

Die zentrale Hypothese lautet: Das mehrstufige Verfahren senkt insbesondere
Profilfragmentierung und fehlerhafte Profilupdates, ohne die Rate korrekter
Wiedererkennungen unvertretbar zu reduzieren.


\section{Datenschutz, Sicherheit und Grenzen}

Lokale Verarbeitung und synthetische IDs reduzieren Risiken, anonymisieren die
Daten jedoch nicht automatisch. Ein ReID-Embedding ist gerade dafür bestimmt,
Beobachtungen wieder zuzuordnen. Daher gelten folgende technische Leitplanken:

\begin{itemize}[leftmargin=7mm]
  \item Zweck und Rechtsgrundlage vor jeder Datenerhebung dokumentieren;
  \item Rohvideos, Snapshots, Events und Embeddings getrennt behandeln und
    jeweils eine Löschfrist definieren;
  \item standardmäßig keine Klarnamen an synthetische IDs binden;
  \item Zugriff, Export, Profil-Merge, Trennung und Löschung auditierbar machen;
  \item Verschlüsselung und Rollenmodell vor Mehrbenutzerbetrieb ergänzen;
  \item Unsicherheit sichtbar halten und Fehlentscheidungen korrigierbar machen;
  \item sensible Zusatzsignale nur als separat freizugebende Module entwickeln.
\end{itemize}

Technisch bleibt OSNet bei starken Kleidungswechseln, Verdeckungen und großen
Perspektivunterschieden begrenzt. Tracking-IDs können fragmentieren; ReID-Scores
sind daten- und domänenabhängig. Die geplante IoT-Anbindung darf diese Grenzen
nicht hinter einer scheinbar eindeutigen Personen-ID verstecken.


\section{Arbeitsplan und Ausblick}

\subsection{Arbeitspakete}

\begin{table}[H]
  \centering
  \small
  \caption{Vorgeschlagener Arbeitsplan für die nächste Projektphase}
  \label{tab:plan}
  \begin{tabularx}{\textwidth}{L{18mm}L{25mm}Y}
    \toprule
    \textbf{Zeitraum} & \textbf{Arbeitspaket} & \textbf{Ergebnis} \\
    \midrule
    Woche 1--2 & Baseline und Governance & Fixierter Datenstand, Testprotokoll, Daten- und Löschkonzept \\
    Woche 3--4 & Unit- und Integrationstests & Tests für Matching-Zonen, Quality Gate, Store-Vertrag und Same-Frame-Ausschluss \\
    Woche 5--6 & Benchmark & ByteTrack/BoT-SORT, YOLO-/OSNet-Varianten und Parametervergleich \\
    Woche 7--8 & Schwierige Szenen & Multi-Person, ähnliche Kleidung, Verdeckung, Wiedereintritt und unbekannte Personen \\
    Woche 9--10 & Multi-Kamera/IoT-Prototyp & Kamera-Registry, Zeitfenster, minimiertes Ereignisschema und ungefährlicher MQTT-Demonstrator \\
    Woche 11--12 & Auswertung & Fehleranalyse, Metriken, Grenzen, reproduzierbarer Ergebnisbericht und Dokumentation \\
    \bottomrule
  \end{tabularx}
\end{table}

\subsection{Mögliche Entwicklungswege}

Nach einer belastbaren Baseline verzweigt sich das Projekt in vier klar
abgrenzbare Wege:

\begin{description}[style=nextline,leftmargin=8mm]
  \item[Weg A -- Smart Home und Assistenz]
    Multi-Kamera-Zonenmodell, lokaler Ereignisbus, Policy Engine und
    datensparsame IoT-Adapter. Dies ist der direkteste Weg zur im Projekt
    motivierten Gebäudeanwendung.
  \item[Weg B -- Genauigkeit und Skalierung]
    Größere Detektoren, Tracker-Tuning, Cross-Camera-ReID, Qdrant-Indizes,
    Profilverwaltung sowie systematisches Sample-Pruning.
  \item[Weg C -- Sportanalyse]
    Teamklassifikation, eigener Ball-/Spielerdetektor, Homographie,
    Meterkoordinaten, Laufwege, Heatmaps und Ballnähe. Die Grundmodule sind
    vorbereitet, aber noch nicht vollständig verdrahtet.
  \item[Weg D -- optionale Sensorfusion]
    Pose, RGB-D, UWB, mmWave oder Wi-Fi-CSI zunächst nur als unabhängige
    Plausibilitätssignale mit eigener Quelle, Konfidenz und Freigabe. Gesicht,
    Stimme, Gangbild und Emotion bleiben besonders sensible Forschungsoptionen.
\end{description}

\subsection{Aktueller Projektstand und erwarteter Beitrag}

\begin{table}[H]
  \centering
  \caption{Abgrenzung von erreichtem Stand und nächster Arbeit}
  \label{tab:status}
  \begin{tabularx}{\textwidth}{L{27mm}Y}
    \toprule
    \textbf{Status} & \textbf{Inhalt} \\
    \midrule
    \statusdone{umgesetzt} & Lokales Windows-Setup, CPU/CUDA-Pfad, YOLO-Tracking, OSNet-ReID, SQLite/Qdrant, Live-UI, Qualitätsfilter, Mehrbild-Puffer, dreistufiges Matching, Profil-Samples, Detail- und Motion-Signale, Evaluationsgrundlage \\
    \statusplan{zu validieren} & Reproduzierbare Multi-Person-Benchmarks, IDF1/HOTA, Threshold-Kalibrierung, Store-Vertragstests, Profilverwaltung und Datenlebenszyklus \\
    \statusplan{geplant} & Multi-Kamera-Architektur, sicherer IoT-Ereignisadapter, Rollen/Rechte, Löschfristen, Football-Vollintegration und optionale Sensorfusion \\
    \bottomrule
  \end{tabularx}
\end{table}

Der erwartete Beitrag ist nicht ein neues neuronales Grundmodell. Er liegt in
der nachvollziehbaren Integration vorhandener Verfahren zu einer lokalen,
konfigurierbaren und evaluierbaren ReID-Pipeline. Besonders relevant sind die
Qualitätssicherung vor dem Speichern, die Wiederverwendung mehrerer guter
Referenzen, die zeitliche Evidenz statt vorschneller Ja/Nein-Entscheidungen und
die klare Trennung zwischen Wahrnehmung, Identitätsannahme und Geräteaktion.
